\documentclass[11pt]{article}
\usepackage[a4paper,margin=27mm]{geometry}
\usepackage{amsmath,amssymb,amsthm,booktabs,array,longtable}
\usepackage[hidelinks]{hyperref}
\newtheorem{theorem}{Theorem}
\newtheorem{proposition}[theorem]{Proposition}
\newtheorem{corollary}[theorem]{Corollary}
\newcommand{\N}{\mathbb N}
\title{Bounded-height classification for Salinas Schmeis's\texorpdfstring{\\$\omega$}{ omega}-chain extension question}
\author{Haobo Ma\\\small ChronoAI Pte Ltd, Singapore
\and Wenlin Zhang\\\small National University of Singapore\\\small ChronoAI Pte Ltd, Singapore}
\date{November 14, 2026}
\begin{document}
\maketitle
\begin{abstract}
Question 2.32 on page 14 of Remi Salinas Schmeis's \emph{General Criterion
for a Poset to Admit a Partition Into $\omega$-Chains}
(arXiv:2610.29056v1, 2026) asks whether every poset has an
$\omega$-chain-partitionable extension that order-embeds into every such
extension. We answer negatively: even a two-element antichain has no such
extension. Two countable extensions force incompatible behavior for upper
bounds of incomparable elements. The same argument gives the connected
three-element fork as a counterexample and, by order duality, the corresponding
negative answer for $\omega^*$-chains. More generally, a nonempty countable
poset of bounded finite height has such a least extension if and only if it
has a greatest element. Every poset has a partitionable extension, even
with $\omega$ replaced by any nonzero ordinal.
\end{abstract}

\section{The question and two extensions}
Salinas Schmeis's universal extension question fails even for a two-element
antichain. We classify the countable posets of bounded finite height for
which the answer is positive, and exhibit two explicit extensions explaining
the smallest obstruction.

An \emph{$\omega$-chain} is a subset with its inherited order isomorphic to
$(\N,\leq)$; here $0\in\N$. Call a poset \emph{admissible} if it has a partition
into $\omega$-chains. An extension of $P$ is a poset containing an induced copy
of $P$, equivalently a poset $Q$ admitting an order embedding $P\hookrightarrow Q$.
An order embedding preserves and reflects the order.
We work with the axiom of choice, as in the source's conventions
\cite[p.~2]{salinas}.

Salinas Schmeis \cite[Question 2.32, p.~14]{salinas} asks whether there is an
admissible extension $Q$ of every $P$ such that
\begin{equation}\label{universal}
 Q\hookrightarrow Q' \qquad\text{for every admissible extension }Q'\text{ of }P.
\end{equation}
The question requires that ``$Q$ is order-isomorphic to a subposet of $Q'$''
\cite[p.~14]{salinas}. It is a leastness property for
embeddability, not simply inclusion-minimality. The embeddings in
\eqref{universal} need not fix the distinguished copy of $P$. We disprove even
this unrestricted version.

The completion analogy motivating \cite[Question 2.32]{salinas} depends on
the target class and the structure preserved. For example, Abela and Chetcuti
\cite[Theorem 3.14]{abela} give a MacNeille extension theorem for join-regular
sublattices when both the ambient complete lattice and the MacNeille completion
satisfy the join-infinite distributive law. Those targets are complete lattices.
In contrast, our first target below has incomparable elements with no common
upper bound.

Let $F=\{r,a,b\}$ have $r<a$, $r<b$ and $a\parallel b$. Define
$R=\{0,1\}\times\N$ by
\begin{equation}\label{rays}
 (i,n)\leq_R(j,m)\quad\Longleftrightarrow\quad
 (i,n)=(0,0)\ \text{or}\ (i=j\text{ and }n\leq m).
\end{equation}
Thus $(0,0)$ is a least element, and above it are two separate rays. This is a
partial order: the only comparisons between different rays start at the least
element; all other comparisons are those of a single copy of $\N$.
The two fibers $\{i\}\times\N$ form an admissible partition. Map
$(r,a,b)$ to $((0,0),(0,1),(1,0))$ to embed $F$ in $R$.

Define $T$ on $\N$ by
\begin{equation}\label{tail}
 x\leq_T y\quad\Longleftrightarrow\quad
 x\leq y\text{ in }\N\ \text{and}\ (x,y)\ne(1,2).
\end{equation}
Reflexivity and antisymmetry follow from the usual order. For transitivity,
if $x\leq_T y\leq_T z$, the only possible failure would be $x=1,z=2$.
Then $y=1$ or $y=2$, and one of the assumed comparisons is already forbidden.
The even and odd integers partition $T$ into $\omega$-chains, with increasing
enumerations $n\mapsto2n$ and $n\mapsto2n+1$. The map
$(r,a,b)\mapsto(0,1,2)$ embeds $F$ in $T$.

Two structural facts give the obstruction. In $R$, any two elements with a
common upper bound are comparable: unless one is the least element, both
must lie on the same ray. In $T$, the only unordered incomparable pair is
$\{1,2\}$, and every strict upper bound of either member is at least $3$ and
therefore an upper bound of both.

\begin{theorem}\label{obstruction}
No poset containing an incomparable pair and having no maximal element
order-embeds into both $R$ and $T$.
\end{theorem}
\begin{proof}
Suppose $x\parallel y$ in $Q$ and $f:Q\hookrightarrow R$,
$g:Q\hookrightarrow T$ are order embeddings. Choose $z>x$, using the absence
of maximal elements. Order reflection makes $g(x),g(y)$ incomparable, so the
structural fact about $T$ gives $g(y)\leq_T g(z)$. Reflecting this comparison
gives $y\leq z$. Consequently $f(z)$ is a common upper bound of $f(x)$ and
$f(y)$ in $R$. Those two elements are comparable, and reflection under $f$
contradicts $x\parallel y$.
\end{proof}

\begin{corollary}\label{counterexamples}
Neither the two-element antichain nor the connected fork $F$ has an admissible
extension satisfying \eqref{universal}. In particular, Question 2.32 is answered
negatively.
\end{corollary}
\begin{proof}
Every element of an admissible poset has a strict successor in its partition
part, so an admissible poset has no maximal element. Both $R$ and $T$ are
admissible extensions of $F$ and of the antichain $\{a,b\}$. A proposed $Q$
contains an incomparable pair and must embed into both $R$ and $T$, contrary
to Theorem~\ref{obstruction}.
\end{proof}

\section{Existence and bounded-height classification}
\begin{proposition}\label{existence}
Every poset $P$ has an admissible extension.
\end{proposition}
\begin{proof}
On $P\times\N$ put
$(p,n)\leq(q,m)$ if $p<q$, or if $p=q$ and $n\leq m$.
For transitivity, if the first coordinates are not all equal, at least one
strict comparison makes the first coordinate strictly increase across the
two comparisons. If they are all equal, use transitivity in $\N$.
Reflexivity is immediate, and antisymmetry follows first in $P$, then in
$\N$. The map $p\mapsto(p,0)$ is an order embedding. The fibers
$\{p\}\times\N$ form the required partition; if $P$ is empty, use the empty
partition.
\end{proof}

We say that $P$ has \emph{bounded finite height} if there is an integer $h$
bounding the number of points in every chain in $P$; for nonempty $P$,
its \emph{height} is the largest such chain size. Countable includes finite.
We use the following elementary partition fact. A nonempty countable set
can index a partition of $\N$ into infinite cofinal subsets: for $m$ indices
use residue classes modulo $m$; for indices $i\in\N$ use
\[
 S_i=\{2^i(2k+1)-1:k\in\N\}.
\]
Each natural number belongs to exactly one $S_i$, by factoring its successor
into a power of two times an odd number.

\begin{theorem}\label{finite}
A nonempty countable poset $P$ of bounded finite height has an admissible extension satisfying
\eqref{universal} if and only if $P$ has a greatest element.
\end{theorem}
\begin{proof}
Let $U=P\oplus\N$, putting every element of $P$ below every element of the
new $\N$-summand. Index a partition of that summand into infinite cofinal
sets $S_p$ by $p\in P$, as above. The sets $\{p\}\cup S_p$ partition $U$
into $\omega$-chains, so $U$ is admissible.

Suppose first that $P$ has greatest element $t$.
Given an admissible extension $e:P\hookrightarrow Q'$, take
the tail strictly above $e(t)$ in the partition part containing $e(t)$.
Map the new $\N$-summand in increasing order onto this tail and use $e$ on
$P$. This is an order embedding: the tail is above every $e(p)$, and no tail
element can be below an $e(p)\leq e(t)$. Thus $U$ satisfies \eqref{universal}.

Conversely, suppose $P$ has no greatest element, and let its height be $h$.
Every point lies below a maximal point, since a strictly ascending sequence
has at most $h$ points. Its set $M$ of maximal elements therefore has at
least two members. Define $V$ by retaining $P$ and adding
a separate ray $\{(m,n):n\in\N\}$ above each $m\in M$. Put
$p<(m,n)$ exactly when $p\leq_P m$, order each ray by $\N$, and put no
comparisons between different rays. These rules define a partial order.
It is admissible: assign each $p\in P$ to one maximal element above it;
then partition each ray into as many infinite cofinal subsequences as there
are points assigned to it, using the countable partition fact, and prepend
one assigned point to each subsequence.
Each maximal element is assigned to itself, so every indexing set is nonempty.

If an admissible $Q$ satisfied \eqref{universal}, it would embed into both
$U$ and $V$. Its embedding into $U$ forces $Q$ to be upward directed.
Here upward directed means that every pair has a common upper bound.
Indeed, above any $x\in Q$ its partition part contains an increasing
$\omega$-sequence. Its image eventually leaves the bounded-height summand $P$ and
exceeds any specified element of $U$. For each $y\in Q$, reflecting such a
comparison gives a common upper bound of $x$ and $y$ in $Q$.

Let $f:Q\hookrightarrow V$. The image of any $\omega$-chain in $Q$ contains
a ray element $(m,n)$, since $P$ has bounded finite height. Directedness
implies that every element of $f(Q)$ lies in
\[
 (\mathord\downarrow_P m)\ \cup\ \{(m,k):k\in\N\}:
\]
a common upper bound with $(m,n)$ must lie on this same ray. Now compose
$P\hookrightarrow Q$ with $f$. The images of the distinct maximal elements
of $P$ are pairwise incomparable. No such image can lie on the ray, whose
elements are comparable with every element of this displayed set. Nor can
any such image be $m$, which is also comparable with every element of the
set. Thus all maximal elements have images strictly below $m$. Every point
of $P$ lies below a maximal element, so the entire image of $P$ is strictly
below $m$. A chain of $h$ points in $P$ would then map to a chain of $h$
points below $m$; appending $m$ gives $h+1$ points in a chain of $P$,
contradicting its height.
\end{proof}

\begin{corollary}\label{finitecase}
A nonempty finite poset has a least admissible extension exactly when it
has a greatest element. The same equivalence holds for every nonempty
countable antichain: the only positive case is a singleton.
The poset with one bottom point and any finite or countably infinite
number of incomparable maximal points at least two is also a counterexample.
\end{corollary}
\begin{proof}
Finite posets and antichains have bounded finite height. Apply
Theorem~\ref{finite}; an antichain has a greatest element only when it is
a singleton.
The last posets have height two and no greatest element.
\end{proof}

The empty poset has the empty admissible extension. Every one-element poset
has a greatest element, whereas the two-element antichain does not. Hence
the two-point counterexample has smallest possible cardinality. For a nonempty
finite chain, $\N$ itself is a least admissible extension: the chain embeds
as an initial finite segment, and every nonempty admissible poset contains
an $\omega$-chain.

\section{Duality and neighboring partition questions}
\begin{corollary}\label{dual}
The two-element antichain has no extension partitionable into
$\omega^*$-chains that embeds into every such extension. The same is true
of the order dual of $F$. Every poset has an extension partitionable into
$\omega^*$-chains. A nonempty countable poset of bounded finite height has a least such extension if
and only if it has a least element.
\end{corollary}
\begin{proof}
Taking order duals converts an $\omega^*$-chain partition into an
$\omega$-chain partition, preserves order embeddings, and converts the
asserted universal property into \eqref{universal}. Apply
Corollary~\ref{counterexamples}, observing that an antichain is self-dual.
The existence assertion and the bounded-height classification follow in the same way
from Proposition~\ref{existence} and Theorem~\ref{finite}, respectively.
\end{proof}

The bounded-height hypothesis cannot simply be dropped: $\N$ itself is
admissible, so it is its own least admissible extension although it has no
greatest element. This is a generalization not asserted in the posing source.

For the neighboring lifting and skeleton questions we record two direct
consequences of the existence construction. A subset is \emph{Noetherian}
if it contains no $\omega$-chain. The \emph{lifting property} of
\cite[Definition 2.26]{salinas} requires that, for every Noetherian subset
$K$, there be an injection from $K$ into the points outside $K$ strictly
above some point of $K$, taking every point strictly upward.

\begin{corollary}\label{lifting}
The extension $E=P\times\N$ in Proposition~\ref{existence} has the lifting
property, with an explicit fiberwise injection for every Noetherian subset.
\end{corollary}
\begin{proof}
For a Noetherian $K\subseteq E$ every section
$K_p=\{n:(p,n)\in K\}$ is finite: an infinite subset of $\N$ has order
type $\omega$. If $K_p$ is nonempty, set $b_p=1+\max K_p$ and send
$(p,n)\in K$ to $(p,b_p+n)$. Its second coordinate exceeds every member
of $K_p$, so it is outside $K$ and strictly above $(p,n)$. Equal images
have equal first coordinates and then equal second coordinates, proving
injectivity. This is also a direct cardinality witness for local shortness
\cite[Definition 2.14]{salinas}.
\end{proof}

For a nonzero ordinal $\alpha$, an \emph{$\alpha$-skeleton}
\cite[Definition 2.33]{salinas} is a disjoint family of inherited
$\alpha$-chains whose union is cofinal: every point lies below a point of
one of the chains.

\begin{corollary}\label{skeleton}
Every poset has an extension partitionable into $\alpha$-chains for each
nonzero ordinal $\alpha$, and hence an extension with an $\alpha$-skeleton.
For each integer $n\geq2$, an $n$-skeleton need not extend to a partition
of the original poset into $n$-chains.
\end{corollary}
\begin{proof}
Order $P\times\alpha$ by $(p,\xi)\leq(q,\eta)$ if $p<q$, or if $p=q$
and $\xi\leq\eta$. The same order-law argument as in
Proposition~\ref{existence} applies. The map $p\mapsto(p,0)$ embeds $P$,
and the fibers partition the extension into $\alpha$-chains. For the second
claim, a chain with $n+1$ points has an $n$-chain containing its greatest
point, hence an $n$-skeleton. Its cardinality is not divisible by $n$, so
it cannot be partitioned into $n$-chains.
\end{proof}

\section{Coverage and verification}
The following table distinguishes answers to the posed questions, consequences
not posed there, and results already supplied by the posing paper. In particular,
the source itself resolves its local-shortness conjecture negatively.
{\small
\begin{longtable}{@{}>{\raggedright\arraybackslash}p{.27\linewidth}>{\raggedright\arraybackslash}p{.66\linewidth}@{}}
\toprule
Source item in \cite{salinas} & Coverage here\\\midrule
\endfirsthead
\toprule
Source item in \cite{salinas} & Coverage here (continued)\\\midrule
\endhead
Question 2.32, p.~14 & Answered: no. Explicit two-point and connected
three-point counterexamples. Proposition~\ref{existence} proves existence
without leastness.\\
Finite and bounded-height cases of Question 2.32 & Not separately posed;
proved: Theorem~\ref{finite} gives the exact countable bounded-height
classification, including all finite posets.\\
Remark 2.29, p.~13 & The source states order duality. Not posed there;
proved: dual obstruction and classification, Corollary~\ref{dual}.\\
Question 2.34, p.~14; Definition 2.33 & Partial answers: Corollary~\ref{skeleton}
gives partitionable extensions for every $\alpha>0$ and shows that an
$n$-skeleton alone is insufficient for every finite $n\geq2$. General
criteria for the original poset are not covered.\\
Question 2.3, p.~3; Theorem 2.27, p.~12 & Prior proof cited: Salinas Schmeis
(2026), Theorem 2.27, gives the general criterion. Not posed separately;
proved: explicit lifting witness for the lexicographic extensions,
Corollary~\ref{lifting}.\\
Question 2.15, p.~7; Conjecture 2.19, p.~8 & Prior proof cited: Salinas Schmeis
(2026), Corollaries 2.23 and 2.25, already gives locally short posets without
partitions. No new general criterion is claimed here.\\
Example 2.10, pp.~5--6; Remark 2.29, p.~13 & Prior proof cited: Salinas
Schmeis (2026), for the uncountable chain and its duality examples.
No new partition theorem for those chains is proved here.\\
Question 2.30, p.~13 & Not covered: existence and values of the reflection
cardinal $\kappa_P$.\\
Question 2.31, p.~14 & Not covered: equivalence of the maximal-chain-family
principle to weak choice principles in ZF.\\
\bottomrule
\end{longtable}}

The pinned
\href{https://github.com/the-omega-institute/trureturing/blob/2b5b552e72596710282b713f6587c7020544535d/D5/S3/ConceptDynamics/DagCompletion/SalinasSchmeisOmegaChainExtensionRefutation.lean}{Lean file for the fork obstruction}
contains the declaration \texttt{result :\ $\neg$\ claim}.
Here \texttt{claim} is the universal extension assertion in
\eqref{universal}, with induced order embeddings and literal partitions into
subsets order-isomorphic to $\N$. Its proof uses the fork, the two orders
\eqref{rays}--\eqref{tail}, and their displayed partitions. This declaration
covers the negative answer via the fork. In the namespace
\begin{quote}\small\ttfamily
D5.S3.ConceptDynamics.DagCompletion.\\
SalinasSchmeisOmegaChainExtensionRefutation
\end{quote}
\noindent
\texttt{OmegaPartitionable} means a literal partition into inherited
$\omega$-chains, and \texttt{IsExtension} means an induced order embedding.
The two-point sharpening, arbitrary-poset extension construction, countable
bounded-height classification, dual statements, explicit lifting injection,
and ordinal-chain extension construction have the written proofs above.
The accompanying \texttt{verify.py} checks the orders,
embeddings and structural implications on explicitly bounded sets; its output
states the bounds and is separate from the proofs for the infinite orders.
It also checks the finite-poset constructions for every transitively closed
subrelation of the natural order on at most five points.

\begin{thebibliography}{2}
\bibitem{salinas} Remi Salinas Schmeis,
\emph{General Criterion for a Poset to Admit a Partition Into
$\omega$-Chains}, arXiv:2610.29056v1, October 21, 2026, 15 pp.
\url{https://arxiv.org/abs/2610.29056v1}.
\bibitem{abela} Kevin Abela and Emmanuel Chetcuti,
\emph{Regular Sublattices, O-Closed Ideals and Dedekind--MacNeille Completions},
arXiv:2609.16280v2, October 21, 2026, 14 pp.
\url{https://arxiv.org/abs/2609.16280v2}.
\end{thebibliography}
\smallskip\noindent\small
\copyright\ 2026 Haobo Ma and Wenlin Zhang.
Licensed under Creative Commons Attribution 4.0 International (CC BY 4.0).
\url{https://creativecommons.org/licenses/by/4.0/}
\end{document}
